% !TEX root = ../main.tex
\chapter*{ABSTRACT}
\addcontentsline{toc}{chapter}{ABSTRACT}

{\fontsize{11}{13.2}\selectfont
Process models produced by conventional \textit{process discovery} algorithms on maintenance data are often complex and hard to read, because activity descriptions are written freely by technicians and therefore generate hundreds of distinct labels for what are essentially similar actions. This study employs \textit{Generative AI}, specifically a \textit{Large Language Model} (LLM), as an annotator that standardizes activity labels on every \textit{event} through taxonomy-based semantic abstraction over twelve classes. Maintenance logs from a single Diesel Power Plant (PLTD) unit were processed into two scenarios, an \textit{as-is event log} using \textit{raw labels} and an \textit{enriched event log} using LLM-annotated labels, each discovered with the \textit{Inductive Miner} and evaluated using four \textit{conformance} metrics. The LLM annotation reduced 112 unique activities into 12 categories, cutting the number of \textit{places} by 61.73\% and \textit{transitions} by 63.87\%. \textit{Conformance} testing showed \textit{generalization} rising sharply from 0.0807 to 0.7079, while \textit{fitness} remained high at 0.9570. The drop in \textit{precision} from 0.8716 to 0.5063 reflects reduced \textit{overfitting} of the \textit{as-is} model rather than a loss of quality. These results demonstrate that LLM-based semantic abstraction can transform a large, tangled process model into one that is compact, generalizable, and operationally meaningful.

\vspace{1cm}

\noindent\textbf{Keywords:} \textit{process mining}, \textit{process discovery}, \textit{event abstraction}, \textit{large language model}, semantic abstraction, \textit{conformance checking}, power plant maintenance
}
